\chapter{SAM and OAM interaction (SOI) of LG Beam}
\label{ch:Theory}
% 連結SHG&DFT
% 2-1
\section{Gauge \add{freedom}}

\add{Although electromagnetic fields are the directly observable physical quantities, structured light beams such as Laguerre--Gaussian (LG) beams are often formulated in terms of the scalar and vector potentials. However, the potentials possess gauge freedom, leading to different forms of the wave equations under different gauge conditions. Importantly, the derivation of LG beams relies on the paraxial approximation, whose implementation is not completely independent of the chosen gauge.}

In this section, we examine how gauge choice \add{affects} the formulation of wave equations. We begin by introducing the wave equations under different gauge conditions, and then analyze the compatibility of the paraxial approximation. This discussion clarifies why certain gauges are more suitable than others for describing paraxial optical fields.

\subsection{Gauge Freedom in Wave Equations}

Maxwell's equations describe how the \textbf{electric field} $\boldsymbol{E}(\boldsymbol{r}, t)$ and the \textbf{magnetic field} $\boldsymbol{B}(\boldsymbol{r}, t)$ evolve in space and time. These equations form the fundamental basis for describing electromagnetic waves, including Laguerre--Gaussian beams. In \textbf{free space}, the charge density $\rho$ and current density $\mathbf{J}$ vanish. In this case, Maxwell's equations reduce to four homogeneous differential equations involving the vacuum permittivity $\epsilon_0$, the vacuum permeability $\mu_0$, and the speed of light in vacuum $c=1/\sqrt{\mu_0\epsilon_0}$. 

Within this electrodynamics framework, Maxwell's equations can be rewritten as two coupled wave equations for the \textbf{scalar and vector potentials}, $\phi(\boldsymbol{r}, t)$ and $\boldsymbol{A}(\boldsymbol{r}, t)$, derived in appendix~\ref{ch:Maxwell} as
\begin{align} 
&\nabla^2\phi(\boldsymbol{r}, t)+\frac{\partial }{\partial t}\left(\boldsymbol{\nabla}\cdot\boldsymbol{A}(\boldsymbol{r}, t)\right)=0, \label{max5} \\
&\nabla^2\boldsymbol{A}(\boldsymbol{r}, t)-\frac{1}{c^2}\frac{\partial^2 }{\partial t^2}\boldsymbol{A}(\boldsymbol{r}, t)-\boldsymbol{\nabla}\left(\boldsymbol{\nabla}\cdot\boldsymbol{A}(\boldsymbol{r}, t)+\frac{1}{c^2}\frac{\partial }{\partial t}\phi(\boldsymbol{r}, t)\right)=0. \label{max6}
\end{align}
These equations contain the same physical information as the original Maxwell equations, but are expressed in terms of the scalar and vector potentials.

Further simplification of Eqs.~(\ref{max5}) and (\ref{max6}) can be achieved by exploiting the gauge freedom of the potentials. Below, we introduce two commonly used gauges: the \textbf{Coulomb gauge} and the \textbf{Lorenz gauge}. Superscripts $C$ and $L$ are used to distinguish quantities defined in the Coulomb and Lorenz gauges, respectively. The gauge conditions are given by
\begin{align} 
\text{Coulomb gauge condition}: \quad &\boldsymbol{\nabla}\cdot\boldsymbol{A}^C(\boldsymbol{r}, t)=0,  \label{Ccondition}\\
\text{Lorenz gauge condition}: \quad &\boldsymbol{\nabla}\cdot\boldsymbol{A}^L(\boldsymbol{r}, t)=-\frac{1}{c^2}\frac{\partial }{\partial t}\phi^L(\boldsymbol{r}, t).
\end{align}
Both gauges simplify the wave equations. In the Coulomb gauge, they reduce to
\begin{align} 
&\nabla^2\phi^C(\boldsymbol{r}, t)=0, \text{and} \label{Cwave1}\\
&\nabla^2\boldsymbol{A}^C(\boldsymbol{r}, t)-\frac{1}{c^2}\frac{\partial^2 }{\partial t^2}\boldsymbol{A}^C(\boldsymbol{r}, t)=\frac{1}{c^2}\frac{\partial }{\partial t}\boldsymbol{\nabla}\phi^C(\boldsymbol{r}, t)=0, \label{Cwave2}
\end{align}
where the last equality follows from Eq.~(\ref{Cwave1}).
In the Lorenz gauge, Eqs.~(\ref{max5}) and (\ref{max6}) become
\begin{align} 
&\nabla^2\phi^L(\boldsymbol{r}, t)-\frac{1}{c^2}\frac{\partial^2 }{\partial t^2}\phi^L(\boldsymbol{r}, t)=0, \text{and} \label{Lwave1}\\
&\nabla^2\boldsymbol{A}^L(\boldsymbol{r}, t)-\frac{1}{c^2}\frac{\partial^2 }{\partial t^2}\boldsymbol{A}^L(\boldsymbol{r}, t)=0. \label{Lwave2}
\end{align}
\add{Here, we introduce the \textbf{monochromatic wave} ansatz, in which the electromagnetic fields and potentials are expressed as 
$\boldsymbol{E}(\boldsymbol{r}, t)=\mathrm{Re}\left[\boldsymbol{E}(\boldsymbol{r}) e^{-i\omega t}\right]$,
$\boldsymbol{A}^C(\boldsymbol{r}, t)=\mathrm{Re}\left[\boldsymbol{A}^C(\boldsymbol{r}) e^{-i\omega t}\right]$,
$\boldsymbol{A}^L(\boldsymbol{r}, t)=\mathrm{Re}\left[\boldsymbol{A}^L(\boldsymbol{r}) e^{-i\omega t}\right]$, and
$\phi^L(\boldsymbol{r}, t)=\mathrm{Re}\left[\phi^L(\boldsymbol{r}) e^{-i\omega t}\right]$,
where $\omega$ is the angular frequency of light. Substituting these expressions into Eq.~(\ref{gaugeF}) from Maxwell's equations and applying the corresponding gauge conditions leads to}
\begin{equation} 
\boldsymbol{E}(\boldsymbol{r})=i\omega\boldsymbol{A}^C(\boldsymbol{r})=
i\omega\left(\boldsymbol{A}^L(\boldsymbol{r})+
\frac{\boldsymbol{\nabla}\left(\boldsymbol{\nabla}\cdot\boldsymbol{A}^L(\boldsymbol{r})\right)}{q_0^2}\right).
\label{LCrelation}
\end{equation}
\add{Here $q_0=\omega/c$ denotes the wavenumber of light in free space. Under the monochromatic-wave ansatz, electric field $\boldsymbol{E}(\boldsymbol{r}, t)$ and magnetic field $\boldsymbol{B}(\boldsymbol{r}, t)$ are related to potentials through Eq.~(\ref{LCrelation}) and Eq.~(\ref{BcrulA})}
%
%
%
%
%
%
\subsection{Incompatibility of Gauge Freedom with the Paraxial Approximation}
In principle, the vector potential formulation of LG beams can be derived starting from either Eq.~(\ref{Cwave2}) or Eq.~(\ref{Lwave2}). \add{However, within Coulomb gauge, the directional constraint impose on the transversely polarized wave is ultimately incompatible with the standard paraxial Helmholtz formulation. For clarity, this incompatibility will be discussed in detail after deriving the paraxial Helmholtz equation within the Lorenz gauge framework~\cite{andrews2012angular}.}

\add{The derivation starts from the scalar Helmholtz equation, which is obtained from} Eq.~(\ref{Lwave2}) by imposing both the monochromatic ($\boldsymbol{A}^L(\boldsymbol{r}, t)=\text{Re}\left[\boldsymbol{A}^L(\boldsymbol{r}) e^{-i\omega t}\right]$) and transverse wave ans\"atze ($\boldsymbol{A}^L(\boldsymbol{r})=\hat{\boldsymbol{\varepsilon}}A^L(\boldsymbol{r})$, where we condsider polarization unit vector $\hat{\boldsymbol{\varepsilon}}=\hat{\boldsymbol{\varepsilon}}_{\sigma}=\hat{\boldsymbol{x}}+i \sigma \hat{\boldsymbol{y}},\,\sigma=\pm 1$ as right/left circularly polarized)(see Sec.~\ref{ch:SHE}):

\begin{equation}
\nabla^2 A^L(\boldsymbol{r}) + q_0^2 A^L(\boldsymbol{r}) = 0.
\label{SH}
\end{equation}
The \textbf{paraxial approximation}~\cite{andrews2012angular, kasap2013optoelectronics, pedrotti2007introduction} assumes that the rays have only small inclinations with respect to the optical axis in geometric optics. In wave optics, this corresponds to an angular spectrum dominated by plane-wave components whose wavevectors are close to the propagation direction. Here, we choose the propagation direction along the $\boldsymbol{z}$ axis, such that the dominant component of the wavevector $\boldsymbol{q}_0$ lies along $\hat{\boldsymbol{z}}$.

\begin{wrapfigure}{l}{0.325\textwidth}
    \centering
    \includegraphics[width=\linewidth, trim=8.3cm 5.6cm 8.3cm 5.8cm, clip]{wave_vector.pdf} %左 下 右 上 
    \caption{\add{Illustration of a dominant component of wavevector $\boldsymbol{q}_0$ lies along $\hat{\boldsymbol{z}}$}}
    \label{fig:wave_vector}
\end{wrapfigure}

Under this approximation, we introduce the \textbf{paraxial wave} ansatz, $A^L(\boldsymbol{r}) = U(\boldsymbol{r}) e^{iq_0 z}$, and substitute it into Eq.~(\ref{SH}), which yields 
\begin{equation}
\nabla^2_\parallel U(\boldsymbol{r})
+ \frac{\partial^2}{\partial z^2} U(\boldsymbol{r})
+ 2 i q_0 \frac{\partial}{\partial z} U(\boldsymbol{r}) = 0.
\label{parahelm}
\end{equation}
Here, $\nabla^2_\parallel = \frac{\partial^2}{\partial x^2} + \frac{\partial^2}{\partial y^2}$ denotes the transverse Laplacian. The paraxial wave ansatz assumes that the \textbf{envelope function} $U(\boldsymbol{r})$ \textbf{varies slowly along the} $\boldsymbol{z}$ \textbf{direction} compared to the rapid phase variation of $e^{i q_0 z}$. Consequently, the paraxial approximation implies the following relations:

\begin{minipage}{0.48\textwidth}
    \begin{equation}
    \left| \frac{\partial^2}{\partial z^2} U(\boldsymbol{r}) \right|
    \ll
    \left| \nabla^2_\parallel U(\boldsymbol{r}) \right|,
    \label{paraxial1}
    \end{equation}
\end{minipage}
\hfill % 自動填滿中間的空白，讓兩個公式漂亮地分開
\begin{minipage}{0.48\textwidth}
    \begin{equation}
    \left| \frac{\partial^2}{\partial z^2} U(\boldsymbol{r}) \right|
    \ll
    q_0 \left| \frac{\partial}{\partial z} U(\boldsymbol{r}) \right|.
    \label{paraxial2}
    \end{equation}
\end{minipage}
\vspace{1.75\baselineskip} % 給下方內文留一點標準行高空隙
Neglecting the second-order derivative with respect to $z$ then leads to the paraxial Helmholtz equation
\begin{equation}
\nabla^2_\parallel U(\boldsymbol{r}) + 2 i q_0 \frac{\partial}{\partial z} U(\boldsymbol{r}) \color{G}{ \approx} 0.
\label{PHE}
\end{equation}

\add{We now clarify why a transversely polarized light ($\hat{\boldsymbol{\varepsilon}}$ confined to the $x$--$y$ plane) is not suitable in this derivation under the Coulomb gauge. For a transversely polarized wave, the Coulomb gauge condition gives}
\begin{align} 
\color{G} 
\boldsymbol{\nabla} \cdot \boldsymbol{A}^C(\boldsymbol{r})
&=
\hat{\boldsymbol{\varepsilon}}_{\sigma} \cdot \boldsymbol{\nabla}_\parallel A^C(\boldsymbol{r}) \\
&=
\left( \frac{\partial}{\partial x}\hat{\boldsymbol{x}} + i \sigma \frac{\partial}{\partial y}\hat{\boldsymbol{y}} \right)A^C(\boldsymbol{r}) = 0
\label{Vio}
\end{align}
\add{where $\boldsymbol{\nabla}_\parallel = \frac{\partial}{\partial x}\hat{\boldsymbol{x}} + \frac{\partial}{\partial y}\hat{\boldsymbol{y}}$. This condition imposes a directional constraint in the transverse plane on the vector potential and results in the transverse Laplaian of the potential become zero.}
\begin{equation}
\color{G}
\nabla^2_\parallel A^C(\boldsymbol{r})
=
\left( \frac{\partial}{\partial x} + i \frac{\partial}{\partial y} \right)
\left( \frac{\partial}{\partial x} - i \frac{\partial}{\partial y} \right)
A^C(\boldsymbol{r})
= 0,
\end{equation}
\add{which in turn is incompatible with the approximation in Eq.~(\ref{paraxial1})~\cite{andrews2012angular}. In contrast, a transversely polarized light do not  have such constraint within Lorenz gauge. For Coulomb gauge, a nonzero longitudinal ($z$) component of the vector potential should be introduced, going beyond the simple transverse polarization. Therefore, the Lorenz gauge provides a more natural and consistent framework for deriving the paraxial Helmholtz equation and Laguerre--Gaussian beam solutions.}
%
%
%
%
%
%
\section{Manifestation of SOI in LG beam}

Within the paraxial approximation, the optical helicity $\sigma$ and the topological charge $\ell$ are directly proportional to the spin and orbital angular momentum of light, respectively~\cite{bliokh2015transverse, barnett1994orbital}. This correspondence provides a convenient framework for analyzing the angular momentum properties of LG beams. 

Moreover, when the vector potential formulation of LG beams are expressed in the Coulomb gauge, an SOI can emerge through the additional longitudinal component of the field, giving rise to terms proportional to $\sigma \ell$ and $\sigma + \ell$. Adopting the Coulomb gauge is further justified by the theory of light–matter interaction, where the vector potential is conventionally formulated within both the Coulomb gauge and the angular spectrum representation. \add{Where in angular spectrum representation can we relate the in-plane wavevector to the vector in reciprocal space, which connect to theory of matter.}

In the following section, we first introduce the LG beams solutions in both real-space and angular spectrum representations. We then perform a gauge transformation to reveal the SOI.
%%%%%%%%%%
\begin{comment} 
can be solved in various coordinate systems and yields different beam solutions, such as the \textbf{Hermite--Gaussian} beam, the \textbf{Laguerre--Gaussian} (LG) beam, and the \textbf{Bessel} beam. 
In cylindrical coordinates $(\rho,\phi,z)$, the paraxial Helmholtz equation takes the form
\begin{equation}
\left(
\frac{1}{\rho}\frac{\partial}{\partial\rho}
+
\frac{\partial^2}{\partial\rho^2}
+
\frac{1}{\rho^2}\frac{\partial^2}{\partial\phi^2}
+
2iq_0\frac{\partial}{\partial z}
\right)
U(\rho,\phi,z)=0 .
\end{equation}
\end{comment}
%%%%%%%%%%

\subsection{Solutions of LG Beams}
\paragraph{Real space representation} 

Under the \add{\textbf{long Rayleigh-range} approximation~\cite{romero2002quantum}}\footnote{\add{For a light beam, the Rayleigh range is the propagation distance from the beam waist to the point where the beam cross-sectional area doubles (see Figure~\ref{fig:Rrange}). It is given by $z_R=\frac{1}{2}q_0W_0^2$, where $W_0$ is the beam waist. When the Rayleigh range is much larger than the characteristic propagation distance of interest, diffraction-induced beam spreading becomes negligible. In this limit, the beam waist can be treated as approximately constant, and the envelope function may be regarded as independent of $z$. This simplification is referred to as the long Rayleigh-range approximation.}}, the LG beam solutions of the paraxial Helmholtz equation, Eq.~(\ref{PHE}), can be expressed in cylindrical coordinates as
\begin{equation}
\color{G}
\boldsymbol{A}^{L}_{\sigma, \ell, p}(\boldsymbol{r})
=
\hat{\boldsymbol{\varepsilon}}_\sigma  A^{L}_{\ell, p}(\boldsymbol{r})
=
\hat{\boldsymbol{\varepsilon}}_\sigma 
\, u_{\ell, p}(\boldsymbol{r}_\parallel)
\, e^{i q_0 z},
\label{LGbeam}
\end{equation}
where the position vector is written as $\boldsymbol{r}=(\rho,\phi,z)=(\boldsymbol{r}_\parallel,z)$. Here $\boldsymbol{r}_\parallel$ denotes the in-plane position vector consisting of the radial coordinate $\rho$ and the azimuthal angle $\phi$. The integer $\ell\in\mathbb{Z}$ is the \textbf{topological charge}, and the integer $p\ge0$ is the radial mode index. The vector $\boldsymbol{q}_0 = q_0 \hat{\boldsymbol{z}}$ denotes the wavevector along the propagation direction. The \add{envelope} function is expressed as
\begin{equation}
\color{G}
u_{\ell, p}(\boldsymbol{r}_\parallel)=f_{\ell, p}(\rho)\,e^{i\ell\phi},
\end{equation}
where the radial function is $f_{\ell, p}(\rho)
=
A^{\rm LG}_0\sqrt{\frac{2p!}{\pi(|\ell|+p)!}}
\,
L_p^{|\ell|}
\left(
\frac{2\rho^2}{W_0^2}
\right)
\left(
\frac{\sqrt{2}\rho}{W_0}
\right)^{|\ell|}
\exp\left(
-\frac{\rho^2}{W_0^2}
\right)$, which is given by~\cite{romero2002quantum, peng2022tailoring, sanchez2024enhanced} and $A^{\rm LG}_0$ is a constant amplitude. Here the generalized Laguerre polynomial is defined as $L_p^{|\ell|}(x)
=
\sum_{m=0}^{p}
(-1)^m
\frac{(|\ell|+p)!}{(p-m)!(|\ell|+m)!m!}
x^m$, with $x=\frac{2\rho^2}{W_0^2}$. The parameter $W_0$ denotes the beam waist of the LG beam, while the cylindrical coordinates are related to the Cartesian coordinates through $\rho=\sqrt{x^2+y^2}$ and $\phi=\tan^{-1}(y/x)$. \add{The geometric meaning of the beam waist and the Rayleigh range is illustrated in Figure~\ref{fig:Rrange}.}

\begin{wrapfigure}{l}{0.455\textwidth}
\centering
\includegraphics[width=\linewidth, trim=1.5cm 3cm 1.5cm 3cm, clip]{Rrange.pdf} %左 下 右 上
\caption{\add{Illustration of the evolution of the beam radius $W(z)$ during propagation. The minimum beam radius $W_0$ occurs at the beam waist, and the Rayleigh range $z_R$ is defined as the distance from the waist at which $W(z_R)=\sqrt{2}W_0$.}}
\label{fig:Rrange}
\end{wrapfigure}

Focusing on the minimum cross-section area of the beam ($z=0$), Figure~\ref{fig:ufig} shows that as $\ell$ and/or $p$ increase, $u_{\ell, p}(\boldsymbol{r}_\parallel)$ develops a broader doughnut-shaped intensity profile, accompanied by an increasing number of radial nodes along the radial direction. The phase varies continuously with the azimuthal angle $\phi$, reflecting the helical wavefront associated with orbital angular momentum. In particular, the phase winds by $2\pi \ell$ around the beam center. The beam waist is chosen as $W_0 = 1.5~\mu\mathrm{m}$, which sets the spatial scale of the beam and determines the radial extent of the intensity distribution. \\

\begin{figure}[H]
\centering
\includegraphics[width=0.95\linewidth, trim=0.2cm 0cm 0.2cm 0.2cm, clip]{ufig.pdf}%左 下 右 上
\caption{Squared magnitude (intensity) and real part distribution of $A^{L}_{\ell, p}(\boldsymbol{r})$ at $z=0$. The green dotted line denotes the beam waist $W_0$. The scale bar indicates the length of 1 $\mathrm{\mu m}$.}
\label{fig:ufig}
\end{figure}

%
%
%

\paragraph{Comparison between the Paraxial and Long-Rayleigh approximations}

{\color{red}
In the derivation, both the paraxial and long Rayleigh-range approximations are employed. Although the two approximations both restrict the propagation of the beam, they impose different physical constraints and therefore cannot be directly compared in terms of the propagation angle alone. The paraxial approximation restricts the propagation angle of the constituent plane waves~\cite{siegman1986lasers}, whereas the long Rayleigh-range approximation restricts the variation of the beam waist during propagation:
\begin{align}
\text{Paraxial approximation}: \quad
&\theta
\equiv
\sin^{-1}\left(\frac{q_{\parallel}}{q_0}\right)
< 0.5~\mathrm{rad}\approx 30^\circ,
\\
\text{Long Rayleigh-range approximation}: \quad
&z \ll z_R,
\end{align}
where $q_{\parallel}=\sqrt{q_x^2+q_y^2}$ is the magnitude of the transverse component of the wavevector, and the Rayleigh range is $z_R=\frac{1}{2}q_0W_0^2$. 

\begin{wrapfigure}{L}{0.45\textwidth}
  \centering
  \includegraphics[width=\linewidth, trim=0cm 0cm 0cm 0cm, clip]{approx_compare.pdf}
  \caption{\color{red}Illustration of the beam-radius evolution and far-field divergence angle $\theta_{\mathrm{div}}$.}
  \label{fig:approx_compare}
\end{wrapfigure}

To facilitate the comparison, the long Rayleigh-range condition can be expressed in terms of the far-field divergence angle~\cite{pedrotti2007introduction}, defined as
\begin{equation}
\theta_{\mathrm{div}}\equiv\lim_{z\rightarrow\infty}\frac{dW(z)}{dz}
=\frac{W_0}{z_R},
\label{theta_div}
\end{equation}
where the beam radius is $W(z)=W_0\sqrt{1+(z/z_R)^2}$. Equation~(\ref{theta_div}) shows that, for a given propagation distance $z$, a smaller far-field divergence angle is formed and allows the beam radius to remain approximately constant over a longer propagation distance (see Figure~\ref{fig:approx_compare}).

\noindent
Although $\theta$ and $\theta_{\mathrm{div}}$ describe different physical quantities, they are related through the propagation of the constituent plane waves. The paraxial approximation therefore also places an upper bound on the far-field divergence angle. Taken together, the two approximation criteria suggest that the characteristic divergence should \textbf{roughly} satisfy
\begin{align}
\text{Paraxial approximation}: \quad
&\theta_{\mathrm{div}}
= \theta_{\mathrm{div}}\ll \frac{W_0}{z_R} \leq  30^\circ,
\\
\text{Long Rayleigh-range approximation}: \quad
&\theta_{\mathrm{div}}\ll \frac{W_0}{z},
\end{align}
where the first condition reflects the paraxial constraint on the relevant angular content, while the second expresses the requirement that the transverse spreading over the propagation distance $z$ remain much smaller than the beam waist.
}

%
%
%
\paragraph{Angular spectrum representation}
To derive the vector potential of LG beams in the angular spectrum representation, the potential is expanded \textbf{in a plane-wave basis} as
\begin{equation}
\boldsymbol{A}(\boldsymbol{r}) = \sum_{\boldsymbol{q}} \boldsymbol{\mathcal{A}}(\boldsymbol{q}) e^{i \boldsymbol{q} \cdot \boldsymbol{r}}, 
\label{planwave}
\end{equation}
where $\boldsymbol{\mathcal{A}}(\boldsymbol{q})$ is the complex amplitude associated with each plane-wave component $e^{i\boldsymbol{q}\cdot\boldsymbol{r}}$. The LG beams in Eq.~(\ref{LGbeam}) can be written in this form through a defined \textbf{two-dimensional Fourier transformation} and its corresponding inverse transformation,
\begin{equation}
u_{\ell, p}(\boldsymbol{r}_\parallel) = \frac{1}{(2\pi)^2} \int {\rm d}^2 \boldsymbol{q}_\parallel\, 
F^{\rm LG}_{\ell, p}(\boldsymbol{q}_\parallel)\, e^{i\boldsymbol{q}_\parallel\cdot \boldsymbol{r}_\parallel},
\quad \longleftrightarrow \quad
F^{\rm LG}_{\ell, p}(\boldsymbol{q}_\parallel) = \int {\rm d}^2 \boldsymbol{r}_\parallel\, 
u_{\ell, p}(\boldsymbol{r}_\parallel)\, e^{-i \boldsymbol{q}_\parallel\cdot \boldsymbol{r}_\parallel}.
\label{eq:FourierPair}
\end{equation}
For applications to two-dimensional materials, the in-plane wavevector $\boldsymbol{q}_\parallel$ can be approximately identified with a reciprocal-space vector. Under this approximation, Eq.~(\ref{eq:FourierPair}) can be discretized, and Eq.~(\ref{planwave}) becomes
\begin{equation} 
\color{G}
\boldsymbol{\mathcal{A}}^{L}_{\sigma, \ell, p}(\boldsymbol{r})
=
\sum_{\boldsymbol{q}_\parallel}
\boldsymbol{\mathcal{A}}^L_{\sigma, \ell, p}(\boldsymbol{q}_\parallel)
e^{i\boldsymbol{q}_\parallel\cdot \boldsymbol{r}_\parallel}
e^{i q_0 z}.
\end{equation}
Here the in-plane wavevector is written as $\boldsymbol{q}_\parallel=(q_\parallel,\phi_{\boldsymbol{q}_\parallel})$. The corresponding complex amplitude can be obtained analytically, as summarized in Table~\ref{tab:my_formulas}.
\begin{table}[H]
\centering
\caption{Real space and angular spectrum representations of LG beams}
\label{tab:my_formulas}

% 縮減回 2 欄結構，確保橫向空間最小化
\begin{tabular}{|c|c|}
\hline
\textbf{Real space} & \textbf{Angular spectrum} \\ \hline

% 第一類別分隔列 
\hline
\multicolumn{2}{|c|}{\textbf{Vector potential}} \\ \hline
\rule[-2.5ex]{0pt}{6.5ex} % 調整支柱高度以配合 \displaystyle 的飽滿尺寸
$\displaystyle \boldsymbol{A}^{L}_{\sigma, \ell, p}(\boldsymbol{r})
=
\hat{\boldsymbol{\varepsilon}}_\sigma A^{L}_{\ell, p}(\boldsymbol{r})
=
\hat{\boldsymbol{\varepsilon}}_\sigma
\, u_{\ell, p}(\boldsymbol{r}_\parallel)
\, e^{i q_0 z}$
&
$\displaystyle \boldsymbol{\mathcal{A}}^L_{\sigma, \ell, p}(\boldsymbol{q}_\parallel)
=
\hat{\boldsymbol{\varepsilon}}_\sigma \mathcal{A}^L_{\ell, p}(\boldsymbol{q}_\parallel)
=
\hat{\boldsymbol{\varepsilon}}_\sigma
\frac{1}{\Omega}
\, F_{|\ell|, p}^{\rm LG}(\boldsymbol{q}_\parallel)$ \\ \hline

% 第二類別分隔列
\hline
\multicolumn{2}{|c|}{\textbf{Envelope function}} \\ \hline
\rule[-2.5ex]{0pt}{6.5ex}
$\displaystyle u_{\ell, p}(\boldsymbol{r}_\parallel)
=
f_{\ell, p}(\rho)e^{i \ell \phi}$
&
$\displaystyle \frac{1}{\Omega}F_{\ell, p}^{\rm LG}(\boldsymbol{q}_\parallel)
=
\tilde{F}_{\left|\ell\right|, p}(q_\parallel)e^{i \ell \phi_{\boldsymbol{q}_{\parallel}}}$ \\ \hline

% 第三類別分隔列 (使用 \displaystyle 強制將換行公式內的字體放大至完全一致)
\hline
\multicolumn{2}{|c|}{\textbf{Solutions at $p=0$}} \\ \hline
\rule[-6ex]{0pt}{13ex} % 微調這列雙行公式的上下留白
$\begin{aligned} 
&\displaystyle f_{\ell, p=0}(\rho) = \\
&\displaystyle \frac{A^{\rm LG}_0}{\sqrt{(|\ell|)!\pi/2}}
\left( \frac{\sqrt{2}\rho}{W_0} \right)^{|\ell|} \exp\left( -\frac{\rho^2}{W_0^2} \right)
\end{aligned}$
&
$\begin{aligned} 
&\displaystyle \tilde{F}_{\left|\ell\right|, p=0}(q_\parallel) = \\
& \displaystyle \frac{A^{\rm LG}_0 W_0^2 \left(-i\right)^{\left|\ell\right|}\pi}{\Omega\sqrt{(|\ell|)!\pi/2}}
 \left( \frac{q_{\parallel} W_0}{\sqrt{2}} \right)^{|\ell|} \exp\left( -\frac{q_{\parallel}^2 W_0^2}{4} \right)
\end{aligned}$ \\ \hline 

\end{tabular}
\end{table}
\noindent
where
$
\tilde{F}_{\left|\ell\right|, p}(q_\parallel)
=
\frac{(-i)^{|\ell|}}{\Omega}
2\pi\,
\mathcal{H}_{|\ell|}\!\left[f_{\ell, p}(\rho)\right],
$
and
$
\mathcal{H}_{|\ell|}\!\left[f_{\ell, p}(\rho)\right]
=
\int_0^\infty {\rm d}\rho\;
\rho\, f_{\ell, p}(\rho)\,
J_{|\ell|}(q_\parallel\rho)
$
is the \textbf{Hankel transform} of order $|\ell|$, where $J_{|\ell|}$ is the \textbf{Bessel function of the first kind}. Here $\Omega$ denotes the area of the two-dimensional system

Figure~\ref{fig:ffig} shows that \add{$\mathcal{A}^L_{\ell, p}(\boldsymbol{q}_\parallel)$} exhibits an intensity profile analogous to that of \add{$A^{L}_{\ell, p}(\boldsymbol{r})$}, reflecting the characteristic doughnut-shaped structure in angular spectrum. This correspondence follows directly from the Fourier transform relation between two representations of the LG beam. The doughnut-shaped distribution in angular spectrum indicates that the dominant in-plane wavevectors lie within a finite range determined by the beam waist $W_0$.
\begin{figure}[H]
\centering
\includegraphics[width=1.0\linewidth]{ffig.pdf}
\caption{Squared magnitude (intensity) and real part distribution of $\mathcal{A}^L_{\ell, p=0}(\boldsymbol{q}_\parallel)$. The scale bar indicates the considered $0.1$ times light cone range for WSe$_2$-monolayers, $Q_c = 8.6~\mu\mathrm{m}^{-1}$.}
\label{fig:ffig}
\end{figure}

%
%
%
%
%
%
% 

\subsection{\add{Longitudinal Component of LG Beam : SOI effect}}

In the previous section, Eq.~(\ref{LCrelation}) was introduced to relate the vector potentials defined in the two gauges to the physical fields. Using this relation, we can perform gauge transformation by rewritten Eq.~(\ref{LCrelation}) as
\begin{equation}
\color{G}
\boldsymbol{A}^{C}_{\sigma, \ell, p}(\boldsymbol{r})
=
\boldsymbol{A}^{L}_{\sigma, \ell, p}(\boldsymbol{r})
+
\frac{\boldsymbol{\nabla}
\left(
\boldsymbol{\nabla}\cdot
\boldsymbol{A}^{L}_{\sigma, \ell, p}(\boldsymbol{r})
\right)}
{q_0^2}.
\label{LCrelation2}
\end{equation}
\add{To explicitly incorporate the SAM of light, the polarization unit vector is defined as
$\hat{\boldsymbol{\varepsilon}}\equiv\hat{\boldsymbol{\varepsilon}}_{\sigma}=\frac{1}{\sqrt{2}}(\hat{\boldsymbol{x}}+i\sigma\hat{\boldsymbol{y}})$,
where $\sigma=\pm1$ denotes the spin of the photon. From now on we omit the superscript $C$ to simplify the notations and the vector potential in real space under the Coulomb gauge is obtained as}
\begin{equation}
\color{G}
\boldsymbol{A}_{\sigma, \ell, p}(\boldsymbol{r})
\equiv 
\boldsymbol{A}^{C}_{\sigma, \ell, p}(\boldsymbol{r})
=\tilde{A}^{\parallel}_{\ell, p}(\boldsymbol{r}) \hat{\boldsymbol{\varepsilon}} + \tilde{A}^{z}_{\sigma, \ell, p}(\boldsymbol{r}) \hat{\boldsymbol{z}}.
\end{equation}
\add{For LG beams in the fundamental radial mode $p=0$} 
\begin{equation} 
\color{G}
\begin{aligned}
\tilde{A}^{\parallel}_{\ell, p}(\boldsymbol{r})
&\approx 
f_{\ell,p=0}(\rho)e^{i\ell \phi}e^{i q_0 z}\\
\tilde{A}^{z}_{\sigma, \ell, p}(\boldsymbol{r})
&\approx
\frac{i}{\sqrt{2}q_0\rho}
\left(
\bigl(|\ell|+\sigma \ell\bigr)-\frac{2\rho^2}{W_0^2}
\right)
f_{\ell,p=0}(\rho)
e^{ij\phi}e^{i q_0 z}
\end{aligned}
\label{AinRinC} % \label 移到這裡，緊接在 \end{aligned} 後面
\end{equation}
\add{where}
\[
\add{J\equiv \sigma+\ell}
\]
\add{is the total angular momentum (TAM).} In general, the second term in Eq.~(\ref{LCrelation2}) introduces corrections to all spatial components of the vector potential. However, under the paraxial approximation, these corrections are suppressed. Therefore, we retain only the longitudinal ($z$) component as the leading correction.

The longitudinal component in Eq.~(\ref{AinRinC}) reveals two distinct types of SOI: a $\sigma\ell$ dependence in the amplitude and a $J$ dependence in the phase. Figure~\ref{fig:AinRinC} shows that the magnitude (absolute value) at $z=0$ of this longitudinal component is significantly weaker than that of the transverse components. The $\sigma\ell$ dependence manifests as an asymmetry between cases with the same $\ell$ but opposite spin $\sigma$.

It is worth noting that the longitudinal term is weighted by a factor proportional to $(q_0 W_0^2)^{-1}$, which is inversely related to the Rayleigh range ($\frac{1}{2} q_0 W_0^2$)~\cite{kasap2013optoelectronics} . This indicates that a smaller beam waist $W_0$ leads to a stronger longitudinal component, corresponding to a shorter Rayleigh range and tighter focusing, as shown in the bottom row of Figure~\ref{fig:AinRinC}. The beam waist $W_0$ is chosen to be $1.5~\mu\mathrm{m}$ and $0.5~\mu\mathrm{m}$, representing a typical paraxial beam and a case at the boundary of the paraxial regime, respectively~\cite{forbes2021relevance, klimov2012mapping, nemoto1990nonparaxial}. This demonstrates that the longitudinal component, although small, remains finite even within the paraxial regime.
\begin{figure}[H]
\centering
\includegraphics[width=1.0\linewidth]{AinRinC.pdf}
\caption{Magnitude (absolute value) of \add{$\boldsymbol{A}_{\sigma, \ell, p=0}(\boldsymbol{r})$} at $z=0$ for different components and spin. The $\sigma\ell$ term leads to asymmetry between cases with the same $\ell$ in the longitudinal component (a wavelength of $532~\mathrm{nm}$ is assumed).}
\label{fig:AinRinC}
\end{figure}

%
%
%
%
As for the $\sigma+\ell$ dependence in the phase term, it is retained in the angular spectrum representation and will be discussed below. By combining Eq.~(\ref{LCrelation2}) with the two-dimensional Fourier transformation defined in Coulomb gauge, vector potential can be written as
\begin{equation}
\color{G}
\boldsymbol{A}_{\sigma, \ell, p}(\boldsymbol{r})
=
\frac{1}{(2\pi)^2}
\int {\rm d}^2 \boldsymbol{q}_\parallel
\,
\boldsymbol{\mathcal{A}}_{\sigma, \ell, p}(\boldsymbol{q_\parallel})
e^{i\boldsymbol{q_\parallel}\cdot \boldsymbol{r}_\parallel}
e^{iq_0 z}.
\label{planwaveinC}
\end{equation}
\add{Since we primarily focus on the fundamental radial mode $p=0$, from now on we omit the subscript $p$ to simplify the notations. The complex amplitude under Coulomb gauge in the angular spectrum representation is then written as
\begin{equation}
\color{G}
\boldsymbol{\mathcal{A}}_{\sigma, \ell}(\boldsymbol{q}_\parallel)
=\boldsymbol{\mathcal{A}}_{\sigma, \ell, p=0}(\boldsymbol{q}_\parallel)
=\tilde{\mathcal{A}}^{\parallel}_{\ell}(\boldsymbol{q}_\parallel) \hat{\boldsymbol{\varepsilon}}_\sigma + \tilde{\mathcal{A}}^{z}_{\sigma, \ell}(\boldsymbol{q}_\parallel) \hat{\boldsymbol{z}}.
\label{AinCinQ}
\end{equation}}
The corresponding complex amplitude in the angular spectrum representation is obtained as
\begin{align}
\color{G}
\boldsymbol{\mathcal{A}}_{\sigma, \ell, p}(\boldsymbol{q_\parallel})
&=
\left[
\hat{\boldsymbol{\varepsilon}}_\sigma
-
\left(
\frac{|\boldsymbol{q}|}{q_0}
\right)^2
(\hat{\boldsymbol{\varepsilon}}_\sigma\cdot\hat{\boldsymbol{q}_\parallel})
\hat{\boldsymbol{q}}
\right]
\mathcal{A}^L_{\ell, p}(\boldsymbol{q_\parallel}) \\
&=
\left[
\hat{\boldsymbol{\varepsilon}}_\sigma
-
\left(
\frac{|\boldsymbol{q}|}{q_0}
\right)^2
\frac{\sin\theta_{\boldsymbol{q}}}{\sqrt{2}}
e^{i\sigma\phi_{\boldsymbol{q}_\parallel}}
\hat{\boldsymbol{q}}
\right]
\mathcal{A}^L_{\ell, p}(\boldsymbol{q_\parallel}) \label{sigma_limit} \\
&\approx
\left[
\hat{\boldsymbol{\varepsilon}}_\sigma
-
\frac{\sin\theta_{\boldsymbol{q}}}{\sqrt{2}}
e^{i\sigma\phi_{\boldsymbol{q}_\parallel}}
\hat{\boldsymbol{z}}
\right]
\mathcal{A}^L_{\ell, p}(\boldsymbol{q_\parallel}),
\label{approx_to_1}
\end{align}
where the final approximation follows from the paraxial condition $\cos\theta_{\boldsymbol{q}} = q_0/|\boldsymbol{q}| \approx 1$. The wavevector unit vector is written as
$
\hat{\boldsymbol{q}}
=
\sin\theta_{\boldsymbol{q}}
\left(
\cos\phi_{\boldsymbol{q_\parallel}}\hat{\boldsymbol{x}}
+
\sin\phi_{\boldsymbol{q_\parallel}}\hat{\boldsymbol{y}}
\right)
+
\cos\theta_{\boldsymbol{q}}\hat{\boldsymbol{z}}
\approx
\hat{\boldsymbol{z}},
$
and 
$
\sin\theta_{\boldsymbol{q}} = \frac{q_\parallel}{\sqrt{q_\parallel^2 + q_0^2}}.
$
Substituting the complex amplitude under Lorenz gauge derived previously, each component of the complex amplitude takes the form
\begin{equation} 
\color{G}
\begin{aligned}
\tilde{\mathcal{A}}^{\parallel}_{\ell}(\boldsymbol{q}_\parallel)
&\equiv 
\tilde{\mathcal{A}}^{\parallel}_{\ell, p=0}(\boldsymbol{q}_\parallel)
\approx 
\tilde{F}_{\left|\ell\right|}(q_\parallel)
e^{i\ell \phi_{\boldsymbol{q}_\parallel}} \\
\tilde{\mathcal{A}}^{z}_{\sigma, \ell}(\boldsymbol{q}_\parallel)
&\equiv 
\tilde{\mathcal{A}}^{z}_{\sigma, \ell, p=0}(\boldsymbol{q}_\parallel)
\approx
-\frac{\sin\theta_{\boldsymbol{q}}}{\sqrt{2}}
\tilde{F}_{\left|\ell\right|}(q_\parallel)
e^{iJ\phi_{\boldsymbol{q}_\parallel}},
\end{aligned}
\label{AinRinQ} 
\end{equation}
\add{where $\tilde{F}_{\left|\ell\right|}(q_\parallel)=\tilde{F}_{\left|\ell\right|, p=0}(q_\parallel)$.} \add{This expression shows that the SOI is encoded in the phase of the longitudinal component through the combined dependence $(\sigma+\ell)$, corresponding to the total angular momentum (TAM). This behavior difference is illustrated in Figure~\ref{fig:A_C_transeverse} and Figure~\ref{fig:A_C_longitudinal}.} However, this phase structure is generally difficult to detext experimentally, as it is averaged out in intensity-based measurements. To access this TAM-dependent phase, interference between two LG beams can be employed, forming vector vortex beams (VVBs), which will be discussed in the following section.
\vspace{-3em}
\begin{figure}[H]
\centering
\includegraphics[width=0.85\linewidth]{A_C_transeverse.pdf}
\caption{Squared magnitude (intensity) and real part distribution of $\tilde{\mathcal{A}}^{\parallel}_{\ell}(\boldsymbol{q}_\parallel)$. The scale bar indicates $0.1$ times the light-cone radius for WSe$_2$ monolayers, $Q_c = 8.6~\mu\mathrm{m}^{-1}$.}
\label{fig:A_C_transeverse}
\end{figure}

\begin{figure}[H]
\centering
\includegraphics[width=0.85\linewidth]{A_C_longitudinal.pdf}
\caption{Squared magnitude (intensity) and real part distribution of $\tilde{\mathcal{A}}^{z}_{\sigma, \ell}(\boldsymbol{q}_\parallel)$. The scale bar indicates $0.1$ times the light-cone radius for WSe$_2$ monolayers, $Q_c = 8.6~\mu\mathrm{m}^{-1}$.}
\label{fig:A_C_longitudinal}
\end{figure}

%
%
%
%
%
\subsection{Vector Vortex Beams(VVBs) : Wave Beating}
Subject to decoherence, the information of TAM encoded in the phase \add{is likely} lost when calculating optical transition rates\add{, as pointed out by Oscar Javier Gomez Sanchez \textit{et al.} in Ref.~\cite{sanchez2024enhanced}}. To retain and control this information, they introduced vector vortex beams (VVBs), which are constructed by superimposing two twisted light beams \add{, which in general, these beams carry angular momentaums of the same magnitude but opposite signs. Thus, VVB carries vetoral polarization that varies across the transverse plane. In contrast, scalar vortex beam, single twisted light beams, carries uniform polarization.} The VVB state is expressed as a superposition through
\begin{equation}
\boldsymbol{\mathcal{A}}_{\sigma, \ell, \sigma', \ell'}(\boldsymbol{q}_\parallel;\alpha,\beta)
=
\cos{\left(\frac{\beta}{2}\right)}
\boldsymbol{\mathcal{A}}_{\sigma, \ell}(\boldsymbol{q}_\parallel)
+
e^{i \alpha}
\sin{\left(\frac{\beta}{2}\right)}
\boldsymbol{\mathcal{A}}_{\sigma', \ell'}(\boldsymbol{q}_\parallel).
\end{equation}
The parameter $\alpha$ controls the relative phase between the two beams, while $\beta$ determines their relative weight in the superposition.  \add{Throughout this thesis}, the radial mode index is set to $p=0$, as the primary focus remains on the interplay between SAM and OAM. $\sigma'$ and $\ell'$ denote the second LG beam that carrying a different combination of angular momentum.
\color{red}
Similar to the previous section, a VVB state contains both transverse and longitudinal components, which can be written as
\begin{equation}
\boldsymbol{\mathcal{A}}_{\sigma, \ell, \sigma', \ell'}(\boldsymbol{q}_\parallel;\alpha,\beta)
= \boldsymbol{\tilde{\mathcal{A}}}_{\sigma, \ell, \sigma', \ell'}^\parallel(\boldsymbol{q}_\parallel;\alpha,\beta)
+ \boldsymbol{\tilde{\mathcal{A}}}_{\sigma, \ell, \sigma', \ell'}^z(\boldsymbol{q}_\parallel;\alpha,\beta),
\end{equation}

\noindent\textbf{Analogy 1D wave beating to longitudinal component:} \\[0.5em]

Here we review wave mixing to obtain an intuitive physical picture of how Vector Vortex Beams (VVBs) will behave. To begin with, let us consider two signals and their additive mixing as:
\begin{align*}
\text{Signal 1}(\phi,t) &= \text{Re}\left[ \tilde{A} e^{i(k\phi-\omega t)} \right], \qquad \text{Signal 2}(\phi,t) = \text{Re}\left[ \tilde{A}' e^{i(k'\phi-\omega t+\alpha)} \right], \quad \text{and} \\
\text{Sum} &\equiv \text{Re}\left[ \tilde{A} e^{i(k\phi-\omega t)} + \tilde{A}' e^{i(k'\phi-\omega t+\alpha)} \right].
\end{align*}
\noindent
Below we show how final wave behave under different condition.

\vspace{3em}
\noindent
\begin{minipage}[t]{0.55\textwidth}
1. Standing wave $\tilde{A} = \tilde{A}' = A$ and $k' = -k$ :
\begin{equation}
\text{Sum} = \text{Re}\Big[ 
\underbrace{2A \cos \left( \frac{\Delta k}{2}\phi + \frac{c}{2} \right)}_{\substack{\text{Envelope}}} 
\underbrace{e^{i\left(- \omega t + \frac{c}{2} \right)}}_{\substack{\text{Carrier}}} 
\Big],
\end{equation}

\noindent
where number of nodes is equal to the difference of wave numbers 
$$
\Delta k \equiv k'-k.
$$
An example for standing wave is shown in Figure~\ref{fig:standing}.

\end{minipage}%
\hfill
\begin{minipage}[t]{0.44\textwidth}
\centering
% 利用 \vtop 將圖片頂端與左側標題第一行對齊
\vtop{\null\hbox{\includegraphics[width=\linewidth]{standing.png}}}
\vspace{0.5em}
{\raggedright
\refstepcounter{figure}\label{fig:standing}%
Figure~\thefigure : Example of a standing wave with $\Delta k = -4$ at $t=0$.\par
}
\end{minipage}

\vspace{3em}
\noindent
\begin{minipage}[t]{0.55\textwidth}
$\text{2. Beating wave} \tilde{A} = \tilde{A}'=A \text{ and } k' \neq k$ :
\vspace{-0.5em}
\begin{equation}
\text{Sum} = \text{Re}\Big[ 
\underbrace{2A \cos \left( \frac{\Delta k}{2}\phi + \frac{\alpha}{2} \right)}_{\text{Envelope}} 
\underbrace{e^{i\left( \bar{k}\phi - \omega t + \frac{\alpha}{2} \right)}}_{\substack{\text{Carrier}}} 
\Big],
\end{equation}

where the carrier start to move along $\phi$ direction (The field is rotating) when 
$$
\bar{k} \equiv \tfrac{k' + k}{2}
$$
is non-zero. An example for standing wave is shown in Figure~\ref{fig:standing2}.
\end{minipage}%
\hfill
\begin{minipage}[t]{0.44\textwidth}
\centering
\vtop{\null\hbox{\includegraphics[width=\linewidth]{standing2.png}}}
\vspace{0.5em}
{\raggedright
\refstepcounter{figure}\label{fig:standing2}%
Figure~\thefigure : Example of a beating wave with $\bar{k} = 9/2$ at $t=0$\par
}
\end{minipage}

\vspace{1em}
\noindent
To analogy to VVB, we consider the amplitude is mismatched but still similar by assuming 
$$
|\tilde{A}|+|\tilde{A}'| \gg |\tilde{A}|-|\tilde{A}'|
$$
3. Real amplitudes $\tilde{A} \neq \tilde{A}' \in \mathbb{R} \text{ and } k' \neq -k$
\begin{equation}
\begin{aligned}
\text{Sum} &= \text{Re} \left\{ \left[ (\tilde{A} + \tilde{A}') \cos \left( \tfrac{\Delta k}{2} \phi + \tfrac{\alpha}{2} \right) - i (\tilde{A} - \tilde{A}') \sin \left( \tfrac{\Delta k}{2} \phi + \tfrac{\alpha}{2} \right) \right] e^{i \left( \bar{k} \phi - \omega t + \tfrac{\alpha}{2} \right)} \right\} \\
&\approx \text{Re} \left\{ \underbrace{(\tilde{A} + \tilde{A}') \cos \left( \tfrac{\Delta k}{2} \phi + \tfrac{\alpha}{2} \right)}_{\text{Envelope}} \underbrace{e^{i \left( \bar{k} \phi - \omega t + \tfrac{\alpha}{2} \right)}}_{\text{Carrier}} \right\}
\end{aligned}
\end{equation}
Under the approximation, the formula is the same but we need to keep in mind that the node becomes non-zero, which corresponding to the term $(\tilde{A} - \tilde{A}')$ is non-zero now. An example is shown in the left panel in Figure~\ref{fig:standing3}

\noindent
4. Complex amplitudes $\tilde{A} \neq \tilde{A}' \in \mathbb{C}$ and $k' \neq -k$:
\begin{equation}
\begin{aligned}
\text{Sum} 
&= \text{Re} \left\{ \left[ (|\tilde{A}| + |\tilde{A}'|) \cos \left( \tfrac{\Delta k}{2}\phi + \tfrac{\alpha}{2} \right) - i (|\tilde{A}| - |\tilde{A}'|) \sin \left( \tfrac{\Delta k}{2}\phi + \tfrac{\alpha}{2} \right) \right] e^{i \left( \bar{k}\phi - \omega t + \tfrac{\alpha}{2} + \tfrac{1}{2}(\text{Arg}[\tilde{A}'] + \text{Arg}[\tilde{A}]) \right)} \right\} \\
&\approx \text{Re} \big\{ \underbrace{(|\tilde{A}| + |\tilde{A}'|) \cos \left( \tfrac{\Delta k}{2}\phi + \tfrac{\alpha'}{2} \right)}_{\text{\scriptsize Envelope}} \underbrace{e^{i \left( \bar{k}\phi - \omega t + \tfrac{\alpha}{2} + \tfrac{1}{2}(\text{Arg}[\tilde{A}'] + \text{Arg}[\tilde{A}]) \right)}}_{\text{\scriptsize Carrier}} \big\}
\end{aligned}
\end{equation}
The relative phase now includes contributions from the complex amplitudes such that 
$$\alpha' \equiv \alpha + \text{Arg}[\tilde{A}'] - \text{Arg}[\tilde{A}].$$
The result shows that the envelope (intensity) is
\begin{equation*}
\text{rotated (counter-clockwise)by } -\frac{\alpha'}{\Delta k}.
\end{equation*}
The rotation (shift) can be seen in the example in the right panel of Figure~\ref{fig:standing3}.

\begin{figure}[H]
\centering
\includegraphics[width=0.9\linewidth]{standing3.png}
\caption{Examples for mismatched amplitudes. The left(right) panel shows the case for real(complex) amplitude, where $-\frac{\alpha'}{\Delta k}=\pi$ for the complex amplitude case.}
\label{fig:standing3}
\end{figure}

Analogy to VVB, the longitudianl component is obtained as
\begin{equation}
\begin{aligned}
\tilde{\boldsymbol{\mathcal{A}}}_{\sigma,\ell,\sigma',\ell'}^{z}(\boldsymbol{q}_{\|}, t; \alpha, \beta) 
&= \left[ \cos \left( \tfrac{\beta}{2} \right) \tfrac{-\sin \theta_{\boldsymbol{q}}}{\sqrt{2}} \tilde{F}_{|\ell|}(q_{\|}) e^{i\left( J\phi_{\boldsymbol{q}_{\|}} - \omega t \right)} + \sin \left( \tfrac{\beta}{2} \right) \tfrac{-\sin \theta_{\boldsymbol{q}}}{\sqrt{2}} \tilde{F}_{|\ell'|}(q_{\|}) e^{i\left( J'\phi_{\boldsymbol{q}_{\|}} - \omega t + \alpha \right)} \right] \hat{\boldsymbol{z}} \\
&\approx \underbrace{G_{|\ell|,|\ell'|}^{+}(q_{\|}; \beta) \cos \left( \frac{\Delta J}{2} \phi_{\boldsymbol{q}_{\|}} + \frac{\alpha_{|\ell|,|\ell'|}}{2} \right)}_{\text{Envelope}} \times \underbrace{e^{i\left( \bar{J}\phi_{\boldsymbol{q}_{\|}} - \omega t + \frac{\alpha}{2} - \frac{\pi}{4}(|\ell'|+|\ell|) \right)}}_{\text{Carrier}} \hat{\boldsymbol{z}},
\end{aligned}
\label{longitudinal}
\end{equation}
where the corresponding indeies are shown in Table~\ref{tab:variable_mapping}
\begin{table}[H]
\centering
\caption{Definition of indeies with corresponding wave beating component in comparison }
\label{tab:variable_mapping}

% 縮減橫向欄間距 (預設為 6pt，降至 3pt 確保橫向緊湊)
\setlength{\tabcolsep}{3pt}
\renewcommand{\arraystretch}{1.1}

\begin{tabular}{|c|c|}
\hline
\rule[-0.6ex]{0pt}{2.8ex}\textbf{Wave Beating} & \textbf{VVB} \\ \hline

% 1. Amplitude (標題極緊湊，公式欄單獨拉高上下間距，括號維持標準尺寸)
\hline
\multicolumn{2}{|c|}{\rule[-0.5ex]{0pt}{2.5ex}\textbf{Amplitude}} \\ \hline
\rule[-2.5ex]{0pt}{6.2ex}
$|\tilde{A}| \pm |\tilde{A}'|$ 
  & $G_{|\ell|,|\ell'|}^{\pm}(q_{\|}; \beta) \equiv \Big| -\tfrac{\sin\theta_q}{\sqrt{2}} \cos\left(\tfrac{\beta}{2}\right) \tilde{F}_{|\ell|}(q_{\|}) \Big| \pm \Big| -\tfrac{\sin\theta_q}{\sqrt{2}} \sin\left(\tfrac{\beta}{2}\right) \tilde{F}_{|\ell'|}(q_{\|}) \Big|$ \\ \hline

% 2. Wavenumber
\hline
\multicolumn{2}{|c|}{\rule[-0.5ex]{0pt}{2.5ex}\textbf{Wavenumber}} \\ \hline
\rule[-1.5ex]{0pt}{4.2ex}
$\Delta k \equiv k' - k, \quad \bar{k} \equiv \tfrac{k' + k}{2}$ 
  & $\Delta J \equiv J' - J, \quad \bar{J} \equiv \tfrac{J' + J}{2}$ \\ \hline

% 3. Relative phase
\hline
\multicolumn{2}{|c|}{\rule[-0.5ex]{0pt}{2.5ex}\textbf{Relative phase}} \\ \hline
\rule[-1.5ex]{0pt}{4.2ex}
$\alpha' \equiv \alpha + \text{Arg}[\tilde{A}'] - \text{Arg}[\tilde{A}]$ 
  & $\alpha_{|\ell|,|\ell'|} \equiv \alpha + \text{Arg}[\tilde{F}_{|\ell'|}(q_{\|})] - \text{Arg}[\tilde{F}_{|\ell|}(q_{\|})] = \alpha - \tfrac{\pi}{2}(|\ell'| - |\ell|)$ \\ \hline

\end{tabular}
\end{table}

\noindent
Without the approximation, $G_{\left|\ell\right|,\left|\ell'\right|}^{+}(q_{\parallel}) \gg G_{\left|\ell\right|,\left|\ell'\right|}^{-}(q_{\parallel})\approx0$, the exact longitudinal component is
\begin{equation}
\begin{aligned}
\tilde{\boldsymbol{\mathcal{A}}}_{\sigma,\ell,\sigma',\ell'}^{\boldsymbol{z}}(\boldsymbol{q}_{\|}, t; \alpha, \beta) 
&= \Big[ G_{|\ell|,|\ell'|}^{+}(q_{\|}; \beta) \cos \left( \tfrac{\Delta J}{2} \phi_{\boldsymbol{q}_{\|}} + \tfrac{\alpha_{|\ell|,|\ell'|}}{2} \right) \\
&\quad - i \underbrace{G_{|\ell|,|\ell'|}^{-}(q_{\|}; \beta)}_{\text{Previously neglected}} \sin \left( \tfrac{\Delta J}{2} \phi_{\boldsymbol{q}_{\|}} + \tfrac{\alpha_{|\ell|,|\ell'|}}{2} \right) \Big] \times e^{i \left( \bar{J}\phi_{\boldsymbol{q}_{\|}} - \omega t + \tfrac{\alpha}{2} - \tfrac{\pi}{4}(|\ell'| + |\ell|) \right)} \hat{\boldsymbol{z}}.
\end{aligned}
\label{longitudinal_exact}
\end{equation}

\noindent
The time-evolving field of $\text{Re}\left[\tilde{\boldsymbol{\mathcal{A}}}_{\sigma,\ell,\sigma',\ell'}^z(\boldsymbol{q}_\parallel,t;\alpha,\beta)\right]$ for different VVB cases can be accessed via the QR code in Figure~\ref{fig:longitudinal_real_alpha_QR}, where Figure~\ref{fig:longitudinal_real} shows only the instantaneous field at $\omega t = 0$. In the animation, the field behaves as:
\begin{equation*}
\begin{aligned}
\text{rotating field : } & \bar{J} \neq 0, \\
\text{standing field : } & \bar{J} = 0.
\end{aligned}
\end{equation*}

\begin{wrapfigure}[6]{r}{0.52\linewidth}
  \centering
  \vspace{-1.2em}
  % 稍微拉大圖片佔比，同時保持緊湊
  \begin{minipage}[c]{0.35\linewidth}
    \centering
    \includegraphics[width=\linewidth, trim=16pt 16pt 16pt 16pt, clip]{longitudinal_real_alpha_QR.png}
  \end{minipage}%
  \hspace{0.2cm}%
  % 讓 Caption 獲得更充裕的橫向排版空間（0.61\linewidth）
  \begin{minipage}[c]{0.61\linewidth}
    \raggedright
    \captionsetup{type=figure, skip=0pt, justification=raggedright}
    \caption{Link to the time evolution animation of $\text{Re}\left[\tilde{\boldsymbol{\mathcal{A}}}_{1,\ell,-1,\ell'}^z(\boldsymbol{q}_\parallel,t;\alpha=0,\beta=\pi/2)\right]$ for Figure~\ref{fig:longitudinal_real}.}
    \label{fig:longitudinal_real_alpha_QR}
  \end{minipage}
  \vspace{-1em}
\end{wrapfigure}

\noindent
Similar to the beating wave, this behavior is depicted by the $\phi_{\boldsymbol{q}_{\|}}$-dependent phase factor in the carrier wave in Eq.~(\ref{longitudinal}).

\clearpage
\begin{figure}[htbp]
  \centering
  \includegraphics[width=0.75\linewidth]{longitudinal_real.png}
  \caption{Following Eq.~(\ref{longitudinal_exact}), the distribution of $\text{Re}\left[\tilde{\boldsymbol{\mathcal{A}}}_{1,\ell,-1,\ell'}^z(\boldsymbol{q}_\parallel,t;\alpha=0,\beta=\pi/2)\right]$ at $\omega t =0$. The animation can be viewed through the QR code in Figure~\ref{fig:longitudinal_real_alpha_QR}.}
  \label{fig:longitudinal_real}
\end{figure}

\noindent
On the other hand, the envelope function is depicted by the intensity (squared absolute value), 
\begin{equation}
\begin{aligned}
\left| \tilde{\boldsymbol{\mathcal{A}}}_{\sigma,\ell,\sigma',\ell'}^z(\boldsymbol{q}_{\|}; \alpha, \beta) \right|^2 
&= \big| G_{|\ell|,|\ell'|}^{+}(q_{\|}; \beta) \cos \big( \tfrac{\Delta J}{2} \phi_{\boldsymbol{q}_{\|}} + \tfrac{\alpha_{|\ell|,|\ell'|}}{2} \big) - i \underbrace{G_{|\ell|,|\ell'|}^{-}(q_{\|}; \beta)}_{\text{neglected}} \sin \big( \tfrac{\Delta J}{2} \phi_{\boldsymbol{q}_{\|}} + \tfrac{\alpha_{|\ell|,|\ell'|}}{2} \big) \big|^2 \\
&\approx \big| \underbrace{ G_{|\ell|,|\ell'|}^{+}(q_{\|}; \beta) \cos \big( \tfrac{\Delta J}{2} \phi_{\boldsymbol{q}_{\|}} + \tfrac{\alpha_{|\ell|,|\ell'|}}{2} \big)}_{\text{Envelope}} \big|^2,
\end{aligned}
\label{longitudinal_envelope}
\end{equation}

\begin{figure}[H]
\centering
\includegraphics[width=0.75\linewidth]{longitudinal.png}
\caption{Following Eq.~(\ref{longitudinal_exact}), the distribution of $\left|\tilde{\boldsymbol{\mathcal{A}}}_{1,\ell,-1,\ell'}^z(\boldsymbol{q}_\parallel;\alpha=0,\beta=\pi/2)\right|^2$.}
\label{fig:longitudinal}
\end{figure}

\noindent
which is shown in Figure~\ref{fig:longitudinal}, where the node structure satisfies :

\begin{equation*}
\text{number of nodes} = \Delta J.
\end{equation*}
This structure is described by the envelope function in Eq.~(\ref{longitudinal}). Note that when the approximation, $G_{\left|\ell\right|,\left|\ell'\right|}^{+}(q_{\parallel}) \gg G_{\left|\ell\right|,\left|\ell'\right|}^{-}(q_{\parallel})\approx0$, is not perfectly satisfied, the nodes of the envelope will not be exactly zero, as can be seen in Eq.~(\ref{longitudinal_exact}).

\noindent
Finally, Figure~\ref{fig:longitudinal_alpha} depicts how the intensity changes with the relative phase $\alpha$ within $\alpha_{|\ell|,|\ell'|}$. The four VVBs shown in the upper row are the same as the upper-row VVBs in Figure~\ref{fig:longitudinal}. The different rows depict the distributions with different $\alpha$ values and show that the node structure is rotated(counter-clockwise) by $-\frac{\alpha_{|\ell|,|\ell'|}}{\Delta J}$. Focusing on the variation with $\alpha$, the node structure is

\begin{equation*}
\text{rotated(counter-clockwise) by } -\frac{\alpha}{\Delta J}
\end{equation*}
However, when $\Delta J =0$, the two twisted light fields have the same wavenumber ($\Delta J =0$), and thus they are coherent and exhibit constructive/destructive interference across the whole $\boldsymbol{q}_{\|}$ plane. This behavior can be seen in Eq.~(\ref{longitudinal_envelope}), where the envelope is independent of $\phi_{\boldsymbol{q}_{\|}}$ and the relative phase $\alpha$ tunes the overall magnitude.

\begin{figure}[H]
\centering
\includegraphics[width=0.75\linewidth]{longitudinal_alpha.png}
\caption{Following Eq.~(\ref{longitudinal_exact}), the distribution of $\left|\tilde{\boldsymbol{\mathcal{A}}}_{1,-1,-1,\ell'}^z(\boldsymbol{q}_\parallel,t;\alpha,\beta=\pi/2)\right|^2$ with variying $\alpha$.}
\label{fig:longitudinal_alpha}
\end{figure}


\clearpage
\noindent\textbf{Analogy 2D wave beating to transverse component:} \\[0.25em]
To discuss the transverse component of VVBs, we need to consider the additional polarization components of light. When superposing two twisted light fields, we can rewrite the different polarizations as
\begin{equation*}
\begin{aligned}
\hat{\boldsymbol{x}} + i \sigma \hat{\boldsymbol{y}} &\xrightarrow{\hspace{3.1cm}} (\hat{\boldsymbol{x}} + \bar{\sigma}\hat{\boldsymbol{y}}) - i \tfrac{\Delta\sigma}{2} \hat{\boldsymbol{y}} \\
\hat{\boldsymbol{x}} + i \sigma' \hat{\boldsymbol{y}} &\xrightarrow{\hspace{3.1cm}} (\hat{\boldsymbol{x}} + \bar{\sigma}\hat{\boldsymbol{y}}) + i \tfrac{\Delta\sigma}{2} \hat{\boldsymbol{y}},
\end{aligned}
\end{equation*}
where the difference in SAM is $\Delta\sigma \equiv \sigma' - \sigma$ and the average SAM is $\bar{\sigma} \equiv \tfrac{\sigma' + \sigma}{2}$. Since we are interested in manipulating the polarization, it is natural to consider light with opposite SAM. In this case, we can observe that for different SAM ($\Delta\sigma=\pm2, \bar{\sigma}=0$), the $\hat{\boldsymbol{x}}$ components of the polarizations remain the same, as in 1D wave beating (additive mixing), while the $\hat{\boldsymbol{y}}$ components only differ in sign (subtractive mixing). The wave beating of the two components is shown below under the approximation
\begin{equation*}
|\tilde{A}| + |\tilde{A}'| \gg |\tilde{A}| - |\tilde{A}'|.
\end{equation*}

\noindent
\textbf{In $\hat{\boldsymbol{x}}$ axis (additive mixing):} 
$\text{Re}\big[ \tilde{A} e^{i(k\phi-\omega t)} + \tilde{A}' e^{i(k'\phi-\omega t+\alpha)} \big] \hat{\boldsymbol{x}}$
\begin{equation*}
\approx \text{Re} \left\{ \big( |\tilde{A}| + |\tilde{A}'| \big) \cos \big( \tfrac{\Delta k}{2} \phi + \tfrac{\alpha'}{2} \big) e^{i \left( \bar{k}\phi - \omega t + \tfrac{\alpha}{2} + \tfrac{1}{2} (\text{Arg}[\tilde{A}'] + \text{Arg}[\tilde{A}]) \right)} \right\} \hat{\boldsymbol{x}}
\end{equation*}

\vspace{0.3em}

\noindent
\textbf{In $\hat{\boldsymbol{y}}$ axis (subtractive mixing if $\bar{\sigma}=0$):} 
$\text{Re}\big[ \textcolor{red}{-} i \tfrac{\Delta \sigma}{2} \tilde{A} e^{i(k\phi-\omega t)} \textcolor{red}{+} i \tfrac{\Delta \sigma}{2} \tilde{A}' e^{i(k'\phi-\omega t+\alpha)} \big] \hat{\boldsymbol{y}}$
\begin{equation*}
\approx \text{Re} \left\{ -\tfrac{\Delta \sigma}{2} \big( |\tilde{A}| + |\tilde{A}'| \big) \sin \big( \tfrac{\Delta k}{2} \phi + \tfrac{\alpha'}{2} \big) e^{i \left( \bar{k}\phi - \omega t + \tfrac{\alpha}{2} + \tfrac{1}{2} (\text{Arg}[\tilde{A}'] + \text{Arg}[\tilde{A}]) \right)} \right\} \hat{\boldsymbol{y}}
\end{equation*}

\vspace{0.3em}

\noindent
\textbf{In 2D : }$\text{Re} \left[ \tilde{A} e^{i(k\phi-\omega t)} \big( \hat{\boldsymbol{x}} \textcolor{red}{-} i \tfrac{\Delta \sigma}{2} \hat{\boldsymbol{y}} \big) + \tilde{A}' e^{i(k'\phi-\omega t+\alpha)} \big( \hat{\boldsymbol{x}} \textcolor{red}{+} i \tfrac{\Delta \sigma}{2} \hat{\boldsymbol{y}} \big) \right]$
\begin{equation*}
\quad \approx \text{Re} \Bigg\{ \underbrace{\big( |\tilde{A}| + |\tilde{A}'| \big) \left[ \cos \big( \tfrac{\Delta k}{2} \phi + \tfrac{\alpha'}{2} \big) \hat{\boldsymbol{x}} - \tfrac{\Delta \sigma}{2} \sin \big( \tfrac{\Delta k}{2} \phi + \tfrac{\alpha'}{2} \big) \hat{\boldsymbol{y}} \right]}_{\text{\small Envelope in 2D (polarization)}} \underbrace{e^{i \left( \bar{k}\phi - \omega t + \tfrac{\alpha}{2} + \tfrac{1}{2} (\text{Arg}[\tilde{A}'] + \text{Arg}[\tilde{A}]) \right)}}_{\text{\small Carrier}} \Bigg\}
\end{equation*}



The transverse component is obtained as:
\begin{align*}
&\tilde{\boldsymbol{\mathcal{A}}}_{\sigma,\ell,\sigma',\ell'}^\parallel(\boldsymbol{q}_\parallel;\alpha,\beta) \\
&= \cos\Big(\tfrac{\beta}{2}\Big) \tilde{F}_{|\ell|}(q_\parallel) e^{i\ell\phi_{\boldsymbol{q}_\parallel}} \hat{\boldsymbol{\varepsilon}}_\sigma + e^{i\alpha} \sin\Big(\tfrac{\beta}{2}\Big) \tilde{F}_{|\ell'|}(q_\parallel) e^{i\ell'\phi_{\boldsymbol{q}_\parallel}} \hat{\boldsymbol{\varepsilon}}_{\sigma'} \\
&= \frac{1}{\sqrt{2}} \left[ \cos\Big(\tfrac{\beta}{2}\Big) \tilde{F}_{|\ell|}(q_\parallel) e^{i\ell\phi_{\boldsymbol{q}_\parallel}} \Big( \hat{\boldsymbol{x}} + i\bar{\sigma}\hat{\boldsymbol{y}} - i\tfrac{\Delta\sigma}{2}\hat{\boldsymbol{y}} \Big) + e^{i\alpha} \sin\Big(\tfrac{\beta}{2}\Big) \tilde{F}_{|\ell'|}(q_\parallel) e^{i\ell'\phi_{\boldsymbol{q}_\parallel}} \Big( \hat{\boldsymbol{x}} + i\bar{\sigma}\hat{\boldsymbol{y}} + i\tfrac{\Delta\sigma}{2}\hat{\boldsymbol{y}} \Big) \right] \\
&= \frac{1}{\sqrt{2}} \left[ G_{|\ell|,|\ell'|}^{+}(q_\parallel;\beta) \cos\Big( \tfrac{\Delta\ell}{2}\phi_{\boldsymbol{q}_\parallel} + \tfrac{\alpha_{|\ell|,|\ell'|}}{2} \Big) - iG_{|\ell|,|\ell'|}^{-}(q_\parallel;\beta) \sin\Big( \tfrac{\Delta\ell}{2}\phi_{\boldsymbol{q}_\parallel} + \tfrac{\alpha_{|\ell|,|\ell'|}}{2} \Big) \right] \\
&\qquad \times e^{i \left[ \bar{\ell}\phi_{\boldsymbol{q}_\parallel} + \tfrac{\alpha_{|\ell|,|\ell'|}}{2} + \operatorname{Arg}(\tilde{F}_{|\ell|}) \right]} \Big( \hat{\boldsymbol{x}} + i\bar{\sigma}\hat{\boldsymbol{y}} \Big) \\
&\quad + \frac{1}{\sqrt{2}} \left[ G_{|\ell|,|\ell'|}^{+}(q_\parallel;\beta) \cos\Big( \tfrac{\Delta\ell}{2}\phi_{\boldsymbol{q}_\parallel} + \tfrac{\alpha_{|\ell|,|\ell'|}+\pi}{2} \Big) - iG_{|\ell|,|\ell'|}^{-}(q_\parallel;\beta) \sin\Big( \tfrac{\Delta\ell}{2}\phi_{\boldsymbol{q}_\parallel} + \tfrac{\alpha_{|\ell|,|\ell'|}+\pi}{2} \Big) \right] \\
&\qquad \times e^{i \left[ \bar{\ell}\phi_{\boldsymbol{q}_\parallel} + \tfrac{\alpha_{|\ell|,|\ell'|}+\pi}{2} + \operatorname{Arg}(\tilde{F}_{|\ell|}) \right]} \Big( -i\tfrac{\Delta\sigma}{2}\hat{\boldsymbol{y}} \Big).
\end{align*}

\noindent
In order to modulate the polarization, we consider $\bar{\sigma}=0$:
\begin{align*}
\tilde{\boldsymbol{\mathcal{A}}}_{\sigma,\ell,\sigma',\ell'}^\parallel(\boldsymbol{q}_\parallel;\alpha,\beta)
\Rightarrow\, & \frac{1}{\sqrt{2}} \Big\{ \underbrace{ G_{|\ell|,|\ell'|}^{+}(q_\parallel;\beta) \left[ \cos\Big( \tfrac{\Delta\ell}{2}\phi_{\boldsymbol{q}_\parallel} + \tfrac{\alpha_{|\ell|,|\ell'|}}{2} \Big)\hat{\boldsymbol{x}} - \tfrac{\Delta\sigma}{2} \sin\Big( \tfrac{\Delta\ell}{2}\phi_{\boldsymbol{q}_\parallel} + \tfrac{\alpha_{|\ell|,|\ell'|}}{2} \Big)\hat{\boldsymbol{y}} \right] }_{\text{anti-node}} \\
&\qquad\quad \underbrace{ + iG_{|\ell|,|\ell'|}^{-}(q_\parallel;\beta) \left[ -\sin\Big( \tfrac{\Delta\ell}{2}\phi_{\boldsymbol{q}_\parallel} + \tfrac{\alpha_{|\ell|,|\ell'|}}{2} \Big)\hat{\boldsymbol{x}} - \tfrac{\Delta\sigma}{2} \cos\Big( \tfrac{\Delta\ell}{2}\phi_{\boldsymbol{q}_\parallel} + \tfrac{\alpha_{|\ell|,|\ell'|}}{2} \Big)\hat{\boldsymbol{y}} \right] }_{\text{node}} \Big\} \\
&\qquad\quad \times \underbrace{ e^{i \left[ \bar{\ell}\phi_{\boldsymbol{q}_\parallel} + \tfrac{\alpha_{|\ell|,|\ell'|}}{2} + \operatorname{Arg}(\tilde{F}_{|\ell|}) \right]} }_{\text{carrier wave}}.
\end{align*}

\noindent
If the spatial overlap between the two beams is sufficient such that
$$
G_{\left|\ell\right|,\left|\ell'\right|}^{+}(q_{\parallel})
\gg
G_{\left|\ell\right|,\left|\ell'\right|}^{-}(q_{\parallel})\approx0,
$$
the term proportional to $G_{\left|\ell\right|,\left|\ell'\right|}^{+}(q_{\parallel})$, corresponding to the anti-node in wave beating, dominates the polarization variation with respect to the azimuthal angle $\phi_{\boldsymbol{q}_\parallel}$. We can therefore approximate the transverse component as
\begin{align}
\tilde{\boldsymbol{\mathcal{A}}}_{\sigma,\ell,\sigma',\ell'}^\parallel(\boldsymbol{q}_\parallel;\alpha,\beta) 
&\approx \frac{1}{\sqrt{2}} G_{|\ell|,|\ell'|}^{+}(q_\parallel;\beta) \left[ \cos\left( \frac{\Delta\ell}{2}\phi_{\boldsymbol{q}_\parallel} + \frac{\alpha_{|\ell|,|\ell'|}}{2} \right)\hat{\boldsymbol{x}} - \frac{\Delta\sigma}{2} \sin\left( \frac{\Delta\ell}{2}\phi_{\boldsymbol{q}_\parallel} + \frac{\alpha_{|\ell|,|\ell'|}}{2} \right)\hat{\boldsymbol{y}} \right] \notag \\
&\quad \times e^{i\left[ \bar{\ell}\phi_{\boldsymbol{q}_\parallel} + \frac{\alpha_{|\ell|,|\ell'|}}{2} + \operatorname{Arg}\left[\tilde{F}_{|\ell|}(q_\parallel)\right] \right]}. \label{eq:envelope_approx}
\end{align}

\noindent
Here, we can observe the polarization is linear. The results for $\sigma=-\sigma'=1, \left(\Delta \sigma=-2 \right)$ are shown in Figure~\ref{fig:transverse}.

\vspace{-1em}
\begin{figure}[H]
\centering
\includegraphics[width=\linewidth]{transverse.png}
\caption{Polarization, instantaneous polarization for the VVB $\tilde{\boldsymbol{\mathcal{A}}}_{1,\ell,-1,\ell'}^\parallel(\boldsymbol{q}_\parallel;\alpha=0,\beta=\pi/2)$, and the corresponding handedness of the elliptical polarization. $\Delta\ell$ is indicated in the bottom right of each case, and those with an odd $\Delta\ell$ are indicated by the red boxes. The notation used in the figure is illustrated in Figure~\ref{fig:transverse_label}.}
\label{fig:transverse}
\end{figure}

\vspace{-1.5em}
\noindent
In Figure~\ref{fig:transverse}, the polarization varies with $\phi_{\boldsymbol{q}_\parallel}$. Analogous to wave beating, where the spatial period of the envelope is $2\pi/|\Delta k|$, the polarization pattern has
\begin{center}
$\text{Azimuthal period} = \frac{2\pi}{|\Delta\ell|}.$
\end{center}
Therefore, although the polarization of each VVB differs as $\phi_{\boldsymbol{q}_\parallel}$ varies, VVBs with the same parity of $|\Delta\ell|$ exhibit the same angular period, while displaying either symmetric or asymmetric polarization patterns depending on whether $|\Delta\ell|$ is even or odd.

\begin{minipage}[t]{0.39\textwidth}
\vspace{0pt}
\centering
\includegraphics[width=0.98\linewidth]{transverse_label.png}
\captionof{figure}{Notation for Figure~\ref{fig:transverse}: dotted line (polarization), orange arrow ($t=0$ instantaneous polarization), triangles (elliptical handedness).}
\label{fig:transverse_label}
\end{minipage}
\hfill
\begin{minipage}[t]{0.56\textwidth}
\vspace{0pt}
\noindent
The elliptical polarization arises from the non-zero $G_{\left|\ell\right|,\left|\ell'\right|}^{-}(q_{\parallel})$, although Eq.~\ref{eq:envelope_approx} still dominates the overall polarization pattern. Furthermore, the relative phase $\alpha_{|\ell|,|\ell'|}$ rotate the polarization across the entire azimuthal angle, where $\alpha_{|\ell|,|\ell'|}=\alpha+\frac{\pi}{2}(|\ell'|-|\ell|)$. The results for $\sigma=-\sigma'=1, \left(\Delta \sigma=-2 \right)$ with varying $\alpha$ are shown in Figure~\ref{fig:transverse_alpha}. This is anology to shift of the envelope function in wave beating.
\end{minipage}

\begin{figure}[H]
\centering
\includegraphics[width=\linewidth]{transverse_alpha.png}
\caption{Polarization, instantaneous polarization for the VVB $\tilde{\boldsymbol{\mathcal{A}}}_{1,-1,-1,\ell'}^\parallel(\boldsymbol{q}_\parallel;\alpha,\beta=\pi/2)$, and the corresponding handedness of the elliptical polarization. The three different rows correspond to different $\alpha$ settings. The notation used in the figure is illustrated in Figure~\ref{fig:transverse_label}.}
\label{fig:transverse_alpha}
\end{figure}

\noindent
In this sense, for a fixed opposite-spin combination ($\Delta \sigma = \pm 2$), the additive mixing of two twisted light beams provides a mechanism to modulate the polarization pattern with a periodic spatial variation along the azimuthal angle $\phi_{\boldsymbol{q}_\parallel}$. In other words, the "winding" of the polarization pattern is determined by the OAM beating of the VVB, governed by $\Delta\ell$. The specific polarization state can be further selected via the relative phase $\alpha$, which rotates the polarization across the entire $\phi_{\boldsymbol{q}_\parallel}$ angle. In addition, the phase of the carrier wave depends on the average OAM---analogous to the phase velocity in wave beating---and strictly alters only the instantaneous polarization.

In conclusion, the superposition of two twisted light beams is analogous to one-dimensional wave beating, consisting of an envelope function and a carrier wave. For envelope function, the azimuthal period is determined by the beating, characterized by the difference in their OAM or TAM for transverse or longitudinal component respectively. For the transverse component, the azimuthal expressed in polarization pattern. For longitudinal component, it is intensity. On the other hand, the instantaneous pattern is futher determined by the average OAM or TAM in the carrier wave for transverse or longitudinal component respectively.

Note that we still need to decide $\Delta \sigma= \pm2$ to determine the polarization, while the beating only decide the azimuthal period. In the next chapter, we take advantage of this beating effect to utilize the selective transition of excitons with finite center-of-mass momentum.


\color{black}





